% GENERATED FILE. Edit canonical lessons, then run npm run generate:books.
% canonical-source-hash: fnv1a64:b3771ce14cb354c3
% canonical-lessons: BN-C229-torch, BN-C229-shikol, BN-C229-korat, BN-C229-phite, BN-C229-abohaoya

\chapter{Torch, Chain, Saw, Ribbon, Weather}
\label{ch:bn-torch-chain-saw-ribbon-weather}
\hlchaptermodality{229}

\begin{chapteropening}
\textbf{By the end of this chapter:} \emph{I can say a torch, a chain, a saw, a ribbon and the weather.}

Complete the last lesson of chapter 229: I can say a torch, a chain, a saw, a ribbon and the weather.
\end{chapteropening}

\section[\texorpdfstring{\bn{টর্চ} (ṭôrch)}{টর্চ (ṭôrch)}]{\bn{টর্চ} (ṭôrch) — a torch}

\label{lesson:BN-C229-torch}



\begin{quote}
Before the new one: say the Bengali for luck, then the Bengali for a decision.
\end{quote}

\subsection*{You'll want to know: \bn{টর্চ}}
\textbf{\bn{টর্চ}} — \emph{ṭôrch} — \textquotedblleft{}a torch\textquotedblright{}.

\begin{cousinweb}
\textbf{\bn{টর্চ}} (\emph{ṭôrch}) is a loan from English \emph{torch}.
\end{cousinweb}

\begin{grammarlens}[title={What you've built}]
One more word for this chapter.
\end{grammarlens}

\subsection*{Your turn}
\begin{itemize}
\raggedright
  \item \emph{Say it:} \emph{ṭôrch}
  \item \emph{Say it:} \emph{ṭôrch}, once more
  \item \emph{Say it:} say \emph{ṭôrch}
  \item \emph{Recall:} say the Bengali for luck, then the Bengali for a decision, then say \emph{ṭôrch} again
\end{itemize}

\begin{tcolorbox}[breakable,colback=teal!4,colframe=teal!35!black,title={Before you move on}]
What does \bn{টর্চ} mean? (\textquotedblleft{}A torch\textquotedblright{}.) Say it once more.
\end{tcolorbox}
\section[\texorpdfstring{\bn{শিকল} (shikôl)}{শিকল (shikôl)}]{\bn{শিকল} (shikôl) — a chain}

\label{lesson:BN-C229-shikol}



\begin{quote}
Before the new one: say the Bengali for a decision, then the Bengali for a torch.
\end{quote}

\subsection*{You'll want to know: \bn{শিকল}}
\textbf{\bn{শিকল}} — \emph{shikôl} — \textquotedblleft{}a chain\textquotedblright{}.

\begin{cousinweb}
\textbf{\bn{শিকল}} (\emph{shikôl}) is a tadbhava, a word inherited from Sanskrit and reshaped by sound change. It comes from \emph{śṛṅkhala}, \textquotedblleft{}chain\textquotedblright{} (CDIAL 12580). CDIAL lists Bengali \emph{sikal}, \textquotedblleft{}chain\textquotedblright{}.
\end{cousinweb}

\begin{grammarlens}[title={What you've built}]
One more word for this chapter.
\end{grammarlens}

\subsection*{Your turn}
\begin{itemize}
\raggedright
  \item \emph{Say it:} \emph{shikôl}
  \item \emph{Say it:} \emph{shikôl}, once more
  \item \emph{Say it:} say \emph{shikôl}
  \item \emph{Recall:} say the Bengali for a decision, then the Bengali for a torch, then say \emph{shikôl} again
\end{itemize}

\begin{tcolorbox}[breakable,colback=teal!4,colframe=teal!35!black,title={Before you move on}]
What does \bn{শিকল} mean? (\textquotedblleft{}A chain\textquotedblright{}.) Say it once more.
\end{tcolorbox}
\section[\texorpdfstring{\bn{করাত} (kôrāt)}{করাত (kôrāt)}]{\bn{করাত} (kôrāt) — a saw}

\label{lesson:BN-C229-korat}



\begin{quote}
Before the new one: say the Bengali for a torch, then the Bengali for a chain.
\end{quote}

\subsection*{You'll want to know: \bn{করাত}}
\textbf{\bn{করাত}} — \emph{kôrāt} — \textquotedblleft{}a saw\textquotedblright{}.

\begin{cousinweb}
\textbf{\bn{করাত}} (\emph{kôrāt}) is a tadbhava, a word inherited from Sanskrit and reshaped by sound change. It comes from \emph{karapattra}, \textquotedblleft{}saw\textquotedblright{} (CDIAL 2795). CDIAL lists Bengali \emph{karāt}, \textquotedblleft{}saw\textquotedblright{}, and notes that the Sanskrit form may be shaped by popular etymology, with a possible Dravidian source.
\end{cousinweb}

\begin{grammarlens}[title={What you've built}]
One more word for this chapter.
\end{grammarlens}

\subsection*{Your turn}
\begin{itemize}
\raggedright
  \item \emph{Say it:} \emph{kôrāt}
  \item \emph{Say it:} \emph{kôrāt}, once more
  \item \emph{Say it:} say \emph{kôrāt}
  \item \emph{Recall:} say the Bengali for a torch, then the Bengali for a chain, then say \emph{kôrāt} again
\end{itemize}

\begin{tcolorbox}[breakable,colback=teal!4,colframe=teal!35!black,title={Before you move on}]
What does \bn{করাত} mean? (\textquotedblleft{}A saw\textquotedblright{}.) Say it once more.
\end{tcolorbox}
\section[\texorpdfstring{\bn{ফিতে} (phite)}{ফিতে (phite)}]{\bn{ফিতে} (phite) — a ribbon, a tape}

\label{lesson:BN-C229-phite}



\begin{quote}
Before the new one: say the Bengali for a chain, then the Bengali for a saw.
\end{quote}

\subsection*{You'll want to know: \bn{ফিতে}}
\textbf{\bn{ফিতে}} — \emph{phite} — \textquotedblleft{}a ribbon, a tape\textquotedblright{}.

\begin{cousinweb}
\textbf{\bn{ফিতে}} (\emph{phite}) is usually traced to Portuguese \emph{fita}, \textquotedblleft{}ribbon\textquotedblright{}.
\end{cousinweb}

\begin{grammarlens}[title={What you've built}]
One more word for this chapter.
\end{grammarlens}

\subsection*{Your turn}
\begin{itemize}
\raggedright
  \item \emph{Say it:} \emph{phite}
  \item \emph{Say it:} \emph{phite}, once more
  \item \emph{Say it:} say \emph{phite}
  \item \emph{Recall:} say the Bengali for a chain, then the Bengali for a saw, then say \emph{phite} again
\end{itemize}

\begin{tcolorbox}[breakable,colback=teal!4,colframe=teal!35!black,title={Before you move on}]
What does \bn{ফিতে} mean? (\textquotedblleft{}A ribbon, a tape\textquotedblright{}.) Say it once more.
\end{tcolorbox}
\section[\texorpdfstring{\bn{আবহাওয়া} (ābôhāoyā)}{আবহাওয়া (ābôhāoyā)}]{\bn{আবহাওয়া} (ābôhāoyā) — the weather}

\label{lesson:BN-C229-abohaoya}



\begin{quote}
Before the new one: say the Bengali for a saw, then the Bengali for a ribbon.
\end{quote}

\subsection*{You'll want to know: \bn{আবহাওয়া}}
\textbf{\bn{আবহাওয়া}} — \emph{ābôhāoyā} — \textquotedblleft{}the weather\textquotedblright{}.

\begin{cousinweb}
\textbf{\bn{আবহাওয়া}} (\emph{ābôhāoyā}) is a loan, usually traced to Persian \emph{āb-o-havā}, \textquotedblleft{}water and air; climate\textquotedblright{}. It is Persian \emph{āb}, \textquotedblleft{}water\textquotedblright{}, + \emph{o}, \textquotedblleft{}and\textquotedblright{}, + \emph{havā}, \textquotedblleft{}air\textquotedblright{}, which is itself Arabic.
\end{cousinweb}

\begin{grammarlens}[title={What you've built}]
One more word for this chapter.
\end{grammarlens}

\subsection*{Your turn}
\begin{itemize}
\raggedright
  \item \emph{Say it:} \emph{ābôhāoyā}
  \item \emph{Say it:} \emph{ābôhāoyā}, once more
  \item \emph{Say it:} say \emph{ābôhāoyā}
  \item \emph{Recall:} say the Bengali for a saw, then the Bengali for a ribbon, then say \emph{ābôhāoyā} again
\end{itemize}

\begin{tcolorbox}[breakable,colback=teal!4,colframe=teal!35!black,title={Before you move on}]
What does \bn{আবহাওয়া} mean? (\textquotedblleft{}The weather\textquotedblright{}.) Say it once more.
\end{tcolorbox}
